% ==============================================================================
% CHAPITRE 1 : CONTEXTE ET PROBLÉMATIQUE
% ==============================================================================
\chapter{Contexte et Problématique}

\begin{infobox}[Objectif du chapitre]
Ce chapitre établit le cadre général de l'étude : le contexte éducatif
numérique, la problématique du traitement des retours étudiants, l'étude
des solutions existantes et la valeur ajoutée de notre approche.
\end{infobox}

% ──────────────────────────────────────────────────────────────────────────────
\section{Contexte général : La Révolution Numérique dans l'Éducation}
% ──────────────────────────────────────────────────────────────────────────────

L'enseignement supérieur traverse une transformation profonde depuis
l'avènement des technologies numériques. La massification de l'éducation
universitaire, conjuguée à l'explosion des plateformes d'apprentissage en
ligne (LMS — \textit{Learning Management Systems}), a multiplié
exponentiellement les interactions numériques entre enseignants et étudiants.

Selon plusieurs études pédagogiques récentes, un étudiant génère en moyenne
plusieurs dizaines de commentaires, évaluations et retours textuels par
semestre académique. À l'échelle d'un département ou d'une université,
cela représente des milliers, voire des dizaines de milliers de textes à
analyser chaque année.

\subsection{L'importance du retour étudiant}

Les retours des étudiants constituent un indicateur fondamental de la
qualité pédagogique. Ils permettent aux enseignants d'identifier les points
faibles de leurs cours, d'adapter leurs méthodes pédagogiques et d'améliorer
l'expérience d'apprentissage. Pour les administrations universitaires, ces
données agrégées fournissent une vision macroscopique de la satisfaction des
apprenants et orientent les décisions d'investissement pédagogique.

Cependant, la valeur réelle de ces données ne peut être exploitée que si
elles sont analysées de manière systématique, objective et rapide. C'est
précisément là que réside le défi technologique que ce projet s'efforce
de relever.

\subsection{Les limites du traitement manuel}

Le traitement manuel des commentaires étudiants présente plusieurs
limitations critiques, résumées dans le tableau \ref{tab:limites_manuel} :

\begin{table}[H]
\centering
\caption{Limitations du traitement manuel des commentaires}
\label{tab:limites_manuel}
\renewcommand{\arraystretch}{1.4}
\begin{tabular}{>{\bfseries}p{3.5cm} p{9cm}}
\toprule
\rowcolor{bleuprincipal!15}
\textbf{Limitation} & \textbf{Impact sur la qualité pédagogique} \\
\midrule
Lenteur & Délai de plusieurs semaines avant exploitation, rendant les insights obsolètes \\
\midrule
\rowcolor{gray!8}
Subjectivité & Interprétation variable selon l'analyste, biais cognitifs inconscients \\
\midrule
Coût élevé & Mobilisation de ressources humaines importantes pour une tâche répétitive \\
\midrule
\rowcolor{gray!8}
Non scalable & Impossible à maintenir avec l'augmentation des effectifs \\
\midrule
Exhaustivité & Tendance à se concentrer sur les cas extrêmes, ignorant les nuances \\
\bottomrule
\end{tabular}
\end{table}

% ──────────────────────────────────────────────────────────────────────────────
\section{Problématique Détaillée}
% ──────────────────────────────────────────────────────────────────────────────

La problématique centrale de ce projet se décline en plusieurs dimensions
complémentaires :

\subsection{Dimension linguistique}

Le français, avec ses 80 formes verbales conjuguées, ses nombreuses
exceptions morphologiques et ses structures syntaxiques complexes, représente
un défi particulier pour l'analyse automatique de sentiments. Les spécificités
suivantes méritent une attention particulière :

\begin{itemize}
  \item \textbf{La négation :} ``Ce cours n'est pas mauvais'' exprime en
        réalité quelque chose de positif, mais une approche naïve basée
        sur la présence du mot ``mauvais'' le classerait négativement.
  \item \textbf{L'intensification :} ``Vraiment excellent'', ``plutôt bien'',
        ``assez décevant'' — les adverbes modificateurs changent fondamentalement
        le niveau de sentiment exprimé.
  \item \textbf{L'ironie et le sarcasme :} ``Super, encore un cours
        incompréhensible...'' — nécessite une compréhension contextuelle
        profonde, impossible pour les approches lexicales.
\end{itemize}

\subsection{Dimension technologique}

Les approches traditionnelles d'analyse de sentiments basées sur des
dictionnaires de polarité (comme SentiWordNet ou VADER) ou des modèles
statistiques simples (SVM, Naive Bayes) présentent des performances
insuffisantes sur le français informel utilisé par les étudiants. Le
besoin d'un modèle profond pré-entraîné sur de larges corpus français
s'impose donc naturellement.

\subsection{Dimension applicative}

Au-delà du modèle seul, le système doit être intégré dans une chaîne
complète : ingestion des données, analyse, stockage, visualisation et
accès depuis des interfaces variées (web, mobile). Cette contrainte
d'intégration impose une architecture distribuée bien conçue.

% ──────────────────────────────────────────────────────────────────────────────
\section{Étude de l'Existant}
% ──────────────────────────────────────────────────────────────────────────────

\subsection{Solutions générales d'analyse de sentiments}

\subsubsection{TextBlob et VADER}

TextBlob et VADER sont des outils d'analyse de sentiments populaires pour
l'anglais. VADER utilise un dictionnaire de polarité augmenté de règles
grammaticales. Bien qu'efficaces pour l'anglais informel, ces outils offrent
des performances médiocres en français (précision inférieure à 60\% sur
des textes académiques).

\subsubsection{Google Cloud Natural Language API}

Google propose une API commerciale d'analyse de sentiments supportant
le français. Elle présente l'avantage d'être prête à l'emploi mais souffre
de plusieurs inconvénients majeurs dans notre contexte : coût par requête,
dépendance à une connexion internet, confidentialité des données étudiantes
et absence de personnalisation au vocabulaire académique.

\subsubsection{Modèles BERT multilingues}

\textbf{mBERT} (Multilingual BERT) et \textbf{XLM-RoBERTa} sont des modèles
Transformer entraînés sur des dizaines de langues simultanément. Bien que
supportant le français, leurs performances restent inférieures à un modèle
monolingue spécialisé, car ils partagent leur capacité entre 100+ langues.

\subsection{CamemBERT : l'état de l'art pour le français}

CamemBERT \cite{martin2020camembert} représente l'état de l'art pour le
Traitement du Langage Naturel en français. Développé conjointement par
Inria, Facebook AI Research et le CNRS, il a été entraîné sur 138 Go de
textes français provenant du corpus OSCAR \cite{ortiz2019oscar}. Ses
performances surpassent systématiquement mBERT sur toutes les tâches NLP
en français.

\begin{table}[H]
\centering
\caption{Comparaison des approches d'analyse de sentiments pour le français}
\label{tab:comparaison_approches}
\renewcommand{\arraystretch}{1.4}
\begin{tabular}{lcccp{3.5cm}}
\toprule
\rowcolor{bleuprincipal!15}
\textbf{Approche} & \textbf{Précision} & \textbf{Gratuit} &
\textbf{Hors-ligne} & \textbf{Personnalisable} \\
\midrule
VADER (anglais) & $<$55\% & Oui & Oui & Limitée \\
\rowcolor{gray!8}
TextBlob French & $\sim$58\% & Oui & Oui & Limitée \\
Google NL API & $\sim$72\% & Non & Non & Non \\
\rowcolor{gray!8}
mBERT & $\sim$70\% & Oui & Oui & Oui \\
XLM-RoBERTa & $\sim$74\% & Oui & Oui & Oui \\
\rowcolor{vertprincipal!15}
\textbf{CamemBERT (notre)} & \textbf{79,41\%} & \textbf{Oui} &
\textbf{Oui} & \textbf{Oui} \\
\bottomrule
\end{tabular}
\end{table}

% ──────────────────────────────────────────────────────────────────────────────
\section{Fondements Théoriques}
% ──────────────────────────────────────────────────────────────────────────────

\subsection{L'architecture Transformer}

Le Transformer \cite{vaswani2017attention} est une architecture de réseau
de neurones introduite en 2017 qui a révolutionné le domaine du NLP. Sa
caractéristique principale est le mécanisme d'\textbf{attention multi-têtes}
(\textit{Multi-Head Self-Attention}) qui permet au modèle de pondérer
l'importance de chaque mot dans une séquence par rapport à tous les autres
mots, capturant ainsi des dépendances à longue distance que les architectures
précédentes (RNN, LSTM) peinent à modéliser.

La formule d'attention est définie par :
\begin{equation}
\text{Attention}(Q, K, V) = \text{softmax}\left(\frac{QK^T}{\sqrt{d_k}}\right) V
\label{eq:attention}
\end{equation}

où $Q$ (Queries), $K$ (Keys) et $V$ (Values) sont des matrices de
projections linéaires, et $d_k$ est la dimension des clés.

\subsection{BERT et le pré-entraînement bidirectionnel}

\acrfull{bert} \cite{devlin2019bert} est un modèle Transformer qui exploite
le contexte bidirectionnel : contrairement aux modèles auto-régressifs comme
GPT qui lisent le texte de gauche à droite, BERT lit simultanément les
contextes gauche et droit de chaque token. Cette bidirectionnalité est
obtenue grâce à deux tâches de pré-entraînement :

\begin{itemize}
  \item \textbf{Masked Language Modeling (MLM) :} 15\% des tokens sont
        masqués aléatoirement, et le modèle doit les prédire.
  \item \textbf{Next Sentence Prediction (NSP) :} le modèle prédit si
        deux phrases sont consécutives dans le texte original.
\end{itemize}

\subsection{CamemBERT}

CamemBERT reprend l'architecture RoBERTa \cite{liu2019roberta} (optimisation
de BERT), mais est entraîné uniquement sur du texte français. Il utilise la
tokenisation \textbf{SentencePiece} \cite{kudo2018sentencepiece} avec un
vocabulaire de 32 000 tokens, permettant une meilleure représentation
des mots français et de leurs formes fléchies.

\subsection{Le fine-tuning par transfert d'apprentissage}

Le \textit{transfer learning} consiste à adapter un modèle pré-entraîné
sur une tâche générale à une tâche spécifique en le ré-entraînant sur un
petit jeu de données annoté. Pour la classification de sentiment, on ajoute
une couche de classification linéaire au-dessus du token \texttt{[CLS]}
de CamemBERT, puis on affine l'ensemble des paramètres :

\begin{equation}
P(\text{classe}_k | \mathbf{x}) = \text{softmax}(W \cdot \mathbf{h}_{CLS} + b)_k
\label{eq:classification}
\end{equation}

où $\mathbf{h}_{CLS}$ est la représentation du token \texttt{[CLS]},
$W$ est la matrice de poids de la couche de classification et $b$ le
biais.

% ──────────────────────────────────────────────────────────────────────────────
\section{Valeur Ajoutée du Projet}
% ──────────────────────────────────────────────────────────────────────────────

Notre système apporte plusieurs contributions concrètes par rapport aux
solutions existantes :

\begin{definitionbox}[Contributions principales]
\begin{enumerate}
  \item \textbf{Spécialisation domaine :} fine-tuning sur des commentaires
        académiques français, améliore la précision sur le domaine cible.
  \item \textbf{Système complet :} de la collecte de données à l'affichage
        mobile, en passant par l'API et le dashboard, le système est
        opérationnel bout en bout.
  \item \textbf{Temps réel :} analyse et retour instantané ($<$ 1 seconde),
        contrairement aux systèmes batch traditionnels.
  \item \textbf{Accessibilité :} application mobile Flutter pour les
        étudiants, dashboard web pour les enseignants.
  \item \textbf{Gratuité et confidentialité :} déployable entièrement
        en local, sans dépendance à des services cloud payants.
\end{enumerate}
\end{definitionbox}

\section*{Conclusion du Chapitre}

Ce chapitre a établi le contexte motivant notre projet, analysé les
limites des approches existantes et positionné CamemBERT comme le choix
technologique le plus pertinent pour l'analyse de sentiments en français
académique. Le chapitre suivant détaille l'analyse fonctionnelle et la
conception architecturale du système.

\cleardoublepage
